%% kapa: exam=1 round=30 year=2019 period=2 form=A subject=accounting subject_ko=회계학 seq=35 number=75
%% source: Q-Net 감정평가사 기출문제 (cst00309) artlSeq=5205205 file=extracted/1008/5205205/2019년 제30회 감정평가사 1차시험 문제지/2019년 제30회 감정평가사 1차시험 문제지/2019년도 제30회 감정평가사 1차 2교시 A형.hwp
%% conversion: kordoc 4.24.51 -> tools/gen1.py (auto); review=proofread-sonnet; flags=none
\begin{question}{75}
20x1년 초 영업을 개시한 (주)감평의 20x1년도와 20x2년도의 생산 및 판매와 관련된 자료는 다음과 같다. \par
{\small\centering
\begin{tabular}{|c|c|c|}
\hline
 & \makecell{20x1년} & \makecell{20x2년} \\
\hline
\makecell{생산량} & \makecell{5,000개} & \makecell{10,000개} \\
\hline
\makecell{판매량} & \makecell{4,000개} & \makecell{10,000개} \\
\hline
\makecell{직접재료원가} & \makecell{￦500,000} & \makecell{￦1,000,000} \\
\hline
\makecell{직접노무원가} & \makecell{￦600,000} & \makecell{￦1,200,000} \\
\hline
\makecell{변동제조간접원가} & \makecell{￦400,000} & \makecell{￦800,000} \\
\hline
\makecell{고정제조간접원가} & \makecell{￦200,000} & \makecell{￦250,000} \\
\hline
\makecell{변동판매관리비} & \makecell{￦200,000} & \makecell{￦400,000} \\
\hline
\makecell{고정판매관리비} & \makecell{￦300,000} & \makecell{￦350,000} \\
\hline
\end{tabular}\par}

(주)감평의 20x2년도 전부원가계산에 의한 영업이익이 ￦100,000일 때, 변동원가계산에 의한 영업이익은? (단, 재공품은 없으며 원가흐름은 선입선출법을 가정한다.) \par
\begin{choices}
  \item ￦85,000
  \item ￦115,000
  \item ￦120,000
  \item ￦135,000
  \item ￦140,000
\end{choices}
\end{question}
\answer{2}
